\ifdefined\gkver\else\def\gkver{student}\fi
\documentclass[\gkver]{gaokaozhenti}
\begin{document}

\chapter{2002年全国理综新课程}
%% paperType: 理综物理部分

\begin{enumerate}
\item
%% number: 17
%% paperName: 2002年全国理综新课程第 17 题 6 分
%% typeId: 1
%% type: 单选题
%% chapter: 光学
%% point: 光的折射
%% method: 无
%% score: 6
%% degree: 750
%% duplicateId: 0
%% body:
一束光线从折射率为 1.5 的玻璃内射向空气，在界面上的入射角为 $45^{\circ}$。下面四个光路图中，正确的是\xzanswer{A}

\fourchoices[answer=A, ispicture=true, h=2.6cm]
{\includesvg{figs/17a}}
{\includesvg{figs/17b}}
{\includesvg{figs/17c}}
{\includesvg{figs/17d}}

%% answer: A
%% memo:
\memoanswer{
光从玻璃射向空气，是从光密介质射向光疏介质，由于玻璃的折射率 $n=1.5$，因此发生全反射的临界角 $\sin C=\frac{1}{n}=\frac{2}{3}<\frac{\sqrt{2}}{2}$，所以 $C<45^{\circ}$，因此图中的光会发生全反射，故 A 正确。
}

\item
%% number: 18
%% paperName: 2002年全国理综新课程第 18 题 6 分
%% typeId: 1
%% type: 单选题
%% chapter: 力物体的平衡
%% point: 力矩平衡
%% method: 无
%% score: 6
%% degree: 550
%% duplicateId: 0
%% body:
图中 AC 为竖直墙面，AB 为均匀横梁，其重为 $G$，处于水平位置。BC 为支撑横梁的轻杆，它与竖直方向成 $\alpha$ 角。A、B、C 三处均用铰链连接。轻杆所承受的力为\xzanswer{D}

\onepicture[h=3.2cm]{figs/18.png}

\fourchoices[answer=D]
{$G\cos\alpha$}
{$\frac{G\cos\alpha}{2}$}
{$\frac{G}{\cos\alpha}$}
{$\frac{G}{2\cos\alpha}$}

%% answer: D
%% memo:
\memoanswer{
由题意可知，B 为活结，BC 为轻杆，因此轻杆 BC 受力一定沿杆方向，则知 BC 杆对横梁 AB 的作用力方向沿 CB 方向向上，如图所示。设 AB 长为 $L$，AB 为匀质横梁，中心在中点处，由三力汇交原理知，铰链对 A 端的作用力方向由 BC 中点指向 A 点，由对称性知，作用力的大小与 $F$ 相等，对 AB，由平衡条件得 $F\cos\alpha=\frac{G}{2}$，解得 $F=\frac{G}{2\cos\alpha}$，故 D 正确。

\onepicture[h=3.4cm]{figs/18_an.jpg}
}

\item
%% number: 19
%% paperName: 2002年全国理综新课程第 19 题 6 分
%% typeId: 1
%% type: 单选题
%% chapter: 交流电
%% point: 交流电
%% method: 无
%% score: 6
%% degree: 650
%% duplicateId: 0
%% body:
在三相交流电源上按星形接法连接相同负载 1、2、3，如图所示，$NN'$ 是中性线。已知负载 1 上的电压为 $220\UV$，电流强度为 $15\UA$。现以 $I$ 表示中性线上的电流，$U$ 表示图中 P、Q 两点之间的电压，则\xzanswer{D}

\onepicture[h=3.0cm]{figs/19.png}

\fourchoices[answer=D]
{$I=15\UA$，$U=440\UV$}
{$I=45\UA$，$U=380\UV$}
{$I=0\UA$，$U=440\UV$}
{$I=0\UA$，$U=380\UV$}

%% answer: D
%% memo:
\memoanswer{
由题意可知，每根火线与地线的电压为 $220\UV$，所以三相交流电线电压（P、Q 之间电压是线电压）为 $380\UV$，而又因为三个完全相同的电阻，所以流过中性线 $NN'$ 的电流为 0，故 D 正确。
}

\item
%% number: 20
%% paperName: 2002年全国理综新课程第 20 题 6 分
%% typeId: 1
%% type: 单选题
%% chapter: 电磁感应
%% point: 双杆问题
%% method: 无
%% score: 6
%% degree: 550
%% duplicateId: 0
%% body:
图中 MN、GH 为平行导轨，AB、CD 为跨在导轨上的两根横杆，导轨和横杆均为导体。有匀强磁场垂直于导轨所在的平面，方向如图。用 $I$ 表示回路中的电流。\xzanswer{C}

\onepicture[h=3.0cm]{figs/20.png}

\fourchoices[answer=C]
{当 AB 不动而 CD 向右滑动时，$I\neq0$ 且沿顺时针方向}
{当 AB 向左、CD 向右滑动且速度大小相等时，$I=0$}
{当 AB、CD 都向右滑动且速度大小相等时，$I=0$}
{当 AB、CD 都向右滑动，且 AB 速度大于 CD 时，$I\neq0$ 且沿逆时针方向}

%% answer: C
%% memo:
\memoanswer{
当 AB 不动而 CD 棒向右运动时，切割产生感应电动势，在整个回路中产生感应电流，其方向逆时针方向，故 A 错误；当 AB 向左、CD 向右滑动且速度大小相等时，此时，穿过电路中的磁通量变大，则整个回路中的感应电流不为零，故 B 错误；若 AB、CD 都向右滑动且速度大小相等时，$I=0$，故 C 正确；当 AB、CD 都向右滑动，且 AB 速度大于 CD 时，则相当于 AB 棒向右切割，因而产生感应电流，$I\neq0$ 且沿顺时针方向，故 D 正确。
}

\end{enumerate}

\end{document}
